\sectionhead{\textbf{About this copy}}

This is a corrected copy of the paper cited above. The text is retyped from the printed paper.
It keeps the paper's wording and the numbering of its results, equations and footnotes, and a
marker such as \origpagemark{$N$} shows where page $N$ of the printed paper begins.

The corrections come from a formalization of the Graph Minors series in the Lean proof
assistant, \href{https://github.com/mahdi-jfri/GraphMinor}{github.com/mahdi-jfri/GraphMinor},
in which the results of this paper are proved. Where the formalization found that a statement
or a step of a proof is wrong as printed, this copy corrects it in place. Where the paper leaves a proof or a step to the reader, this
copy adds it, also in blue. Text \del{struck
through in red} is removed and text \fix{in blue} is inserted. Each change has a note, marked
\textcolor{fixcolor}{\sffamily F1}, \textcolor{fixcolor}{\sffamily F2}, \dots, that gives the
reason. The file
\href{https://github.com/mahdi-jfri/GraphMinor/blob/main/corrections/GM15.md}{corrections/GM15.md}
of the formalization describes each correction in full. Each note, and each item in the list below, begins with the kind of the change. A \kindname{misprint} is a wrong symbol, index or word. A \kindname{statement} is a result or a claim of the text that is false as printed. A \kindname{degenerate} is a statement that fails only in a trivial case, such as an empty graph or a parameter equal to~$0$. A \kindname{gap} is a step of a proof that does not follow as written. A \kindname{missing} is a proof that the paper omits, leaves to the reader or only sketches, and that this copy adds.

The corrections and added proofs, in the order of the pages:
\begin{enumerate}\setlength{\itemsep}{1pt}
\item \kindtag{fix:1.1-proof-theta}\hyperref[fix:1.1-proof-theta]{Proof of (1.1), page 116}: $\theta_{k+1}$ and
$\lambda_{k+1}$ become $\theta_{\kappa+1}$ and $\lambda_{\kappa+1}$.
\item \kindtag{fix:1.1-proof-lambda}\hyperref[fix:1.1-proof-lambda]{Proof of (1.1), page 116}: $\lambda_{k+1}$ becomes
$\lambda_{\kappa+1}$.
\item \kindtag{fix:2.1-proof-empty}\hyperref[fix:2.1-proof-empty]{Proof of (2.1), page 116}: the case
$\bar{\Omega} = \emptyset$ is set aside first.
\item \kindtag{fix:2.1-claim8}\hyperref[fix:2.1-claim8]{Claim (8) in the proof of (2.1), page 118}:
$N(P_i \cup B)$ becomes $N(P_i \cap B)$.
\item \kindtag{fix:2.1-proof-R}\hyperref[fix:2.1-proof-R]{Proof of (2.1), page 119}: $V(A \cup B)$ becomes
$V(A \cap B)$.
\item \kindtag{fix:3.4}\hyperref[fix:3.4]{Statement of (3.4), page 122}: the paths $P_i$ are mutually disjoint
and meet $H$ only in their ends.
\item \kindtag{fix:3.4-iv-step-added}\hyperref[fix:3.4-iv-step-added]{Proof of (3.4), page 123}: the proof that (3.2)(iv)
holds is added.
\item \kindtag{fix:3.4-6.1-step-added}\hyperref[fix:3.4-6.1-step-added]{Proof of (3.4), page 123}: the application of Graph
Minors XII, (6.1), is given in full.
\item \kindtag{fix:3.4-proof-indices}\hyperref[fix:3.4-proof-indices]{Proof of (3.4), page 123}: the indices of $e_1, e_2$ and
of $v_1, v_2$ in the last sentence are corrected.
\item \kindtag{fix:3.5}\hyperref[fix:3.5]{Statement of (3.5), page 123}: a hypothesis on the $H$-normal lines
from $bd(\Lambda_1)$ to $bd(\Lambda_2)$ is added.
\item \kindtag{fix:4-ins}\hyperref[fix:4-ins]{Section 4, page 124}: ``radial graph of $G$'' becomes ``radial
drawing of $H$''.
\item \kindtag{fix:5.2}\hyperref[fix:5.2]{Statement of (5.2), page 129}: $\mu \ge 0$ becomes $\mu \ge 1$.
\item \kindtag{fix:5.2-claim7}\hyperref[fix:5.2-claim7]{Claim (7) in the proof of (5.2), page 131}: $\Lambda_1$ becomes
$Z_1$.
\item \kindtag{fix:5.2-claim8}\hyperref[fix:5.2-claim8]{Claim (8) in the proof of (5.2), page 131}: $\Lambda_2$ becomes
$Z_2$.
\item \kindtag{fix:6.2}\hyperref[fix:6.2]{Statement of (6.2), page 134}: $2\tau_i$ becomes $2\tau_1$.
\item \kindtag{fix:7-pi}\hyperref[fix:7-pi]{Section 7, page 135}: in the definition of $\pi$,
$4 + \mu + \theta_{k+1}/2$ becomes $7 + \mu + \theta_{k+1}/2$.
\item \kindtag{fix:7-tau2}\hyperref[fix:7-tau2]{Section 7, page 138}: $\tau_2 = 3 + \theta'/2$ becomes
$\tau_2 = 6 + \theta'/2$.
\item \kindtag{fix:7.7-claim8}\hyperref[fix:7.7-claim8]{Claim (8) in the proof of (7.7), page 142}: $j \le i$ becomes
$j \le k$.
\item \kindtag{fix:7.11-claim1}\hyperref[fix:7.11-claim1]{Claim (1) in the proof of (7.11), page 143}: $2\tau - 6$
becomes $2\tau - 12$.
\end{enumerate}
